\section{計算主張とその監査の範囲}
\label{sec:computational}

この作業中の草稿における分類は、解析的な議論と以下の6つの計算主張からの
論理的帰結である。これらの計算主張は、現時点で利用できる証拠とは切り離して、
数学的な主張として述べる。
その名称は、保管された計算の位置づけを格上げするものではない。
いくつかのソースパッケージは
\texttt{DIAGNOSTIC\_ONLY / NOT\_EVIDENCE}という指定を保持している。ある結果を用いるという
プロジェクト上の決定は、それ自体では独立な数値的証明にならない。リポジトリは
公開されており、固定された出典識別子は参考文献に記してある。
必要なすべての包囲について独立に実装した再実行があるとは、
本改訂では主張しない。

本節を通じて、$E$は参照座標でのエネルギー
$\widehat E_\lambda$の略記であり、空間微分はすべて参照座標で
とる。次のようにおく。
\[
 \mathcal D=\{(\rho,z):\rho,z\geq0,\ \rho^2+z^2\leq1\},\qquad
 H=\frac{E_\rho}{\rho},\qquad V=\frac{E_z}{z}.
\]
$\rho=0$または$z=0$となる内点では、それぞれの商を
対応する2階導関数によって拡張する。2階以上の空間微分に
関する不等式はすべて、たとえその箱が境界と交わる場合でも、
記載した箱と開球との交わりの上での主張である。
境界では、1階導関数は全勾配の連続な
内側トレースを意味する。このとき$H$は$\rho>0$のところでのみ、
その動径方向のトレースを$\rho$で割ったものとして用い、$V$は
$z>0$のところでのみ用いる。本節ではHessianの境界への拡張は仮定しない。
全勾配のトレース条件は、境界上の制約付き
臨界点条件とは別のものである。

\subsection{有限の言明}

\begin{computationalassertion}[有限扁平域における横断方向と軸上の符号]
\label{comp:c1}
$1/32\leq\lambda\leq1$と内点$(\rho,z)\in\mathcal D$に対して
\[
 H(\rho,z;\lambda)>0.
\]
$\rho>0$である境界点での動径方向のトレースの商についても、
同じ狭義の符号が成り立つ。
$0<z\leq1$に対して、軸上の1階導関数は次を満たす。
\[
 \begin{array}{ll}
 E_z(0,z;\lambda)<0,&1/32\leq\lambda\leq2/5,\\
 E_z(0,z;\lambda)>0,&33/50\leq\lambda\leq1.
 \end{array}
\]
軸上の線分$0\leq z\leq1/8$では、より詳しく、次が成り立つ。
\[
 \begin{array}{ll}
 E_{zz}(0,z;\lambda)<0,&1/32\leq\lambda\leq2/5,\\
 E_{zz}(0,z;\lambda)>0,&33/50\leq\lambda\leq1.
 \end{array}
\]
対応する1階導関数の符号は$z=0$から積分することで
従う。$1/8\leq z\leq1$では$E_z$を直接包囲する。
\end{computationalassertion}

保管された有限被覆は、$H$そのものか$X=E_\rho$のいずれかを用いる。
$X$の葉は、その箱全体で$\rho>0$であり、かつ
$X$の下端が正であるときに限り採用する。したがって、それが$H$について
含意することには、明示された分母の条件が付く。箱を空として捨てるのは、
$\mathcal D$との交わりが空であると示された後に限る。
パラメータのスラブの重なりは意図したものである。$2/5$は
上向きの二進端点を通して被覆され、上側の族は
$21/32<33/50$から始まる。
$[21/32,1]$において、$H$、$z\geq1/8$での軸上の$E_z$、
および$z\leq1/8$での軸上の$E_{zz}$について記録された下端は、それぞれ次のとおりである。
\[
 \frac{11239578}{2^{36}},\qquad
 \frac{28921220}{2^{36}},\qquad
 \frac{34030169745}{2^{36}}.
\]
計算主張~\ref{comp:c1}は、この有限な扁平側の帯における軸外の臨界点を
すべて排除する。残りの窓における軸上の個数調査は、
以下の別個のスカラーの言明である。$0<\lambda<1/32$での一様な
横断方向の符号については何も主張しない。
後者の範囲は、薄い扁平の場合の解析的議論によって扱う\citep{data-oblate-full}。

\begin{computationalassertion}[扁平側における軸上のスカラー個数調査]
\label{comp:c2}
次のようにおく。
\[
 g(t,\lambda)=E_z(0,t;\lambda),\quad
 h_o(\lambda)=g_t(0,\lambda),\quad
 c_3(\lambda)=\frac16 g_{ttt}(0,\lambda),\quad
 b(\lambda)=g(1,\lambda).
\]
次の言明が成り立つ。
\begin{enumerate}
 \item $[1/4,1]$上で$h_o'(\lambda)>0$であり、
 $h_o(2/5)<0<h_o(83/200)$である。したがって、その唯一の零点$c_o$は次を満たす。
 \[
  \frac25<c_o<\frac{83}{200}.
 \]
 さらに、$[2/5,83/200]$全体で$c_3(\lambda)<0$である。
 \item 極でのトレースは$[2/5,33/50]$に唯一の零点$b_o$をもち、
 その零点の前では負、後では正であって、次が成り立つ。
 \[
  \frac{674802}{2^{20}}<b_o<\frac{674811}{2^{20}}.
 \]
 記録された端点での不等式は次のとおりである。
 \[
  b(674802/2^{20})\leq-\frac{175746}{2^{36}},\qquad
  b(674811/2^{20})\geq\frac{153930}{2^{36}}.
 \]
 \item $[2/5,33/50]$上で、$g$の正の内部零点は次に従う。
 \[
  \#\{t\in(0,1):g(t,\lambda)=0\}
  =\begin{cases}
    1,&c_o<\lambda<b_o,\\
    0,&2/5\leq\lambda\leq c_o\ \text{or}\ b_o\leq\lambda\leq33/50.
   \end{cases}
 \]
 そのような零でない内部の根のすべてにおいて$g_t<0$である。
 $\lambda=b_o$では、端点$t=1$が$(0,1]$における唯一の正の零点
 である。これは内部零点には数えない。
\end{enumerate}
\end{computationalassertion}

スカラーの根の個数は、指定した符号ゲートから組み立てるのであって、
描画した枝の標本から得るのではない。下半分では
$g(t,\lambda)=t\Phi(t^2,\lambda)$と書く。
$[2/5,9/20]\times[0,1/2]$上で、保管されたエネルギーの4階導関数の
ゲート$g_{ttt}<0$は、解析的な移送恒等式と合わせて、
$\partial_\tau\Phi<0$を与える。すると、$t=1/2$と中心での符号付きの値が
下半分の根を決定する。$[2/5,9/20]$上の上半分の被覆は
160個のパラメータのスラブをもち、根のない48個のスラブ、
端点の符号と$g_t$が負の帯状領域をもつ78個のスラブ、および
孤立した根の管をもつ34個のスラブに分かれる。$[9/20,5/8]$上では、採用された159個の
パラメータ区間が、左側での正の値、根の管における$g_t<0$、
および右側での負の値または負の極でのトレースをもつ。
極の手前で終わる管は124個、極で終わる管は35個である。
前者の124個に対する残りの外側領域は別途被覆される。
$[5/8,33/50]\times[31/32,1]$上では、ゲートは$g_t<0$、
$g(31/32,\lambda)>0$、および符号付きの極でのトレースである。この記述は、
どの不等式が存在、一意性、排除を与えるかを特定する。
完全な閉区間での組み立てとその出典記録は、
軸に関する参照資料に保管されている\citep{data-oblate-axis}。

C2(3)の論理的な使い方を明示するため、$\lambda_x$を
$(699/1600,281/640)$における$g(1/2,\lambda)$の唯一の零点
とする（C1aゲート）。以下は、固定された軸上の記録に基づく
\emph{ゲートから根の個数への組み立て}である
\citep{data-oblate-axis}。ここで、$t>0$に対して$\Phi(t^2,\lambda)=g(t,\lambda)/t$、
$B_o(\lambda)=g(1,\lambda)$であり、「管」は
端点の符号が逆で$g_t<0$である区間を表し、$N_+$は
$(0,1)$内の零点のみを数える。各項目は指定した符号ゲートを条件とするものであり、
それらの区間実装を独立に検証するものではない。

\begin{center}
\begingroup\footnotesize
\setlength{\tabcolsep}{3pt}\renewcommand{\arraystretch}{1.12}
\begin{tabularx}{\linewidth}{@{}p{29mm}p{43mm}Xr@{}}
\toprule
$\lambda$ & 軸の下側$0<t\le1/2$ & 軸の上側$1/2<t<1$と外側の排除 & $N_+$\\
\midrule
$[2/5,\cO)$ & $\Phi(0)<0$、$\Phi_\tau<0$：根なし & Gap-1 L：$g<0$ & $0$\\
$\{\cO\}$ & $\Phi(0)=0$、$\Phi_\tau<0$：正の根なし & Gap-1 L：$g<0$ & $0$\\
$(\cO,83/200]$ & $\Phi(0)>0>\Phi(1/4)$、狭義単調減少：根は1つ & Gap-1 L：$g<0$ & $1$\\
$(83/200,\lambda_x)$ & C1a/A1、$\Phi(0)>0>\Phi(1/4)$：根は1つ & Gap-1 X：上側の区間は排除 & $1$\\
$\{\lambda_x\}$ & $\Phi(1/4)=0$：$t=1/2$のみ & X：$t>1/2$はすべて排除 & $1$\\
$(\lambda_x,9/20)$ & $\Phi(1/4)>0$：下側に根なし & X/R：単調な上側の根が1つ、外側領域は負 & $1$\\
$[9/20,5/8]$ & C1c：$g>0$ & C1b：採用された管159個、外側領域124個とGap-2は排除 & $1$\\
$(5/8,\bO)$ & C1d：$31/32$まで$g>0$ & 帯：$g_t<0$、$g(31/32)>0>B_o$ & $1$\\
$\{\bO\}$ & C1d：$31/32$まで$g>0$ & 帯：$g_t<0$、零点は$t=1$のみ & $0$\\
$(\bO,33/50]$ & C1d：$31/32$まで$g>0$ & 帯：$g_t<0$、$B_o>0$：上側に根なし & $0$\\
\bottomrule
\end{tabularx}
\endgroup
\end{center}
端点$83/200$は第3行に割り当てる。Gap-1のスラブとしては
この値はLではなくXに属し、整合する符号検査を
与える。$703/1600$におけるX/Rの接合部は共有された1つの認証済みの
根であり、2つの根ではない。値$9/20$はC1bに属し、$5/8$は
C1bの行に属する。$5/8$におけるC1d/帯の検査は、その同じ根を
1回だけ数える。切れ目$t=31/32$は記帳上のものであり、一方$t=1$は境界上の
零点であって$N_+$には決して含めない。固定されたソースには、
Lが48個、Xが78個、Rが34個のパラメータのスラブと、クランプなし124個、
クランプあり35個のC1bの管も記録されている。本ZIPには別個に読める
C2の証拠対応表が含まれる。C2のソースコード全体の再実行は別の課題として残る。

以前の中心の計算では、160ビットのArbを用い、続いて書き写した核の192ビットでの再実行を行った\citep{data-center-axis}。別個の端点検査のために、$s=\sin t$、$m=\cos t$、$D_0=s^2+\lambda^2 m^2$、$\theta=\arctan((\lambda-1/\lambda)ms)$と書く。軸上では、符号付き角は$\arctan(\lambda(m-z)/s)-\arctan(m/(\lambda s))$である。重み$1-zm$を付けたその二乗を$z=0$で2回微分し、鏡映を用いると次を得る。
\[
 h_o(\lambda)=\int_0^{\pi/2}
 \frac{s\,[2\lambda^2s^2+4\lambda(1-\lambda^2)m s^3\theta]}{D_0^2}\,dt .
\]
表~\ref{tab:critical}のより細かい区間の端点において、複素ボール積分は左側で上界$-186781876279625303/2^{96}$を、右側で下界$10585972787056429/2^{96}$を与える。C2の単調性と合わせると、これは同じ零点を分離する。C2のより広い区間は、その3次係数の検査も支える。

\begin{computationalassertion}[球に近い帯]
\label{comp:c3}
$1\leq\lambda\leq2$と$(\rho,z)\in\mathcal D$に対して
\begin{align*}
 E_z(\rho,z;\lambda)&\geq\frac{16189}{2^{36}}
       &&(1/4\leq z\leq1),\\
 E_{zz}(\rho,z;\lambda)&\geq\frac{66330426}{2^{36}}
       &&(0\leq z\leq1/4).
\end{align*}
$z=0$上では$0\leq\rho<1$全体で$E_{yyyy}<0$である。境界の赤道上では
\[
 H_b(\lambda):=E_\rho(1,0;\lambda)
 \geq\frac{32725738}{2^{36}}>0.
\]
\end{computationalassertion}

鏡映により$E_z(\rho,0)=0$であるから、最初の2つの不等式から、
$z\ne0$であれば常に$V\geq16189/2^{36}$が従う。
赤道全体での符号と正の境界値により、中心以外の赤道上の点はすべて排除される。以前の15葉による赤道上の検査は、その代わりに、正の半径において$H\geq1383839082/2^{36}$または$E_{\rho\rho}\leq-4138135153/2^{36}$のいずれかを包囲していた。これは別個の一方向交差の議論を与える。
ソースの訂正E2は、球面展開の最初の2項
$4/3-8r^2/9$が厳密な公式であるとした以前の主張を撤回する。
計算主張~\ref{comp:c3}の不等式はいずれも、その打ち切りを用いない\citep{data-near-sphere}。

\begin{computationalassertion}[扁長側の有限な窓]
\label{comp:c4}
$2\leq\lambda\leq5$と$(\rho,z)\in\mathcal D$に対して
\begin{align*}
 E_{zz}&\geq\frac{1503705}{2^{33}}
          &&(0\leq z\leq1/4),\\
 E_z&\geq\frac{3341881}{2^{36}}
          &&(1/4\leq z\leq1).
\end{align*}
赤道上の$0\leq\rho<1$に対して$E_{yyyy}\leq-109964349/2^{36}$である。$H$の正の半径の内部零点のそれぞれにおいて、別個の形状微分の入力が$H_\lambda<0$を与える。
赤道上の$0\leq\rho\leq1/8$に対して、
$E_{yyyy}\leq-5355762265/2^{35}$である。ここで$y=q_2$である。
中心係数$H_c(\lambda)=H(0,0;\lambda)$は次を満たす。
\[
 H_c'(\lambda)\leq-\frac{1090349669}{2^{34}}
 \qquad(2\leq\lambda\leq5).
\]
境界係数は次を満たす。
\[
 \begin{array}{ll}
 H_b'(\lambda)\leq-6439538363/2^{36},&2\leq\lambda\leq9/4,\\
 H_b(\lambda)\leq-11303125/2^{36},&9/4\leq\lambda\leq5.
 \end{array}
\]
境界と中心の唯一の零点をそれぞれ$b_p$、$c_p$とすると、
これらは次によって包囲される。
\[
 b_p\in\left(\frac{34651360}{2^{24}},\frac{34651369}{2^{24}}\right),
 \qquad
 c_p\in\left(\frac{79261996}{2^{24}},\frac{79262005}{2^{24}}\right).
\]
表示した左端点と右端点において、上下界はそれぞれ次のとおりである。
\[
 \begin{array}{c|cc}
 &\text{left endpoint, lower bound}&\text{right endpoint, upper bound}\\ \hline
 H_b&6485992/2^{47}&-5189437/2^{47}\\
 H_c&3088214/2^{47}&-2555623/2^{47}.
 \end{array}
\]
$[2,9/4]\times[7/8,1)\times\{z=0\}$上で、
$E_{\rho\rho}\leq-28849345377/2^{36}$が境界の赤道に至るまで
一様に成り立つ。$\rho=1$では1階導関数のトレースのみを
用いる。
\end{computationalassertion}

最初の2つの不等式は、赤道の外で$V\geq3341881/2^{36}$を
与える。4階導関数の上界は、式~\eqref{eq:fourth_derivative}により、すべての$0<\rho<1$に対して$H_\rho<0$を与える。
別個のパラメータ方向の被覆は、$H>0$、
$H<0$、$H_\lambda<0$という選択肢をもつ。
したがって、すべての根は両方の狭義の導関数不等式を満たし、
陰関数の微分により
$\rho_*'(\lambda)=-H_\lambda/H_\rho<0$を得る。
4次のゲートは退化した中心を制御し、端点の
符号は2つの遷移の位置を定める\citep{data-window}。

\begin{computationalassertion}[扁長側の裾の有限部分]
\label{comp:c5}
$l=\sqrt{\lambda^2-1}$とし、式~\eqref{eq:pt-three-coefficient-densities}の1次元積分
$N_c,N_e,N_p$を用いる。
次のようにおく。
\[
 A_c(l)=3\pi/4-N_c(l),\quad
 A_e(l)=3\pi/4-N_e(l),\quad
 A_p(l)=\pi/2-N_p(l).
\]
アンカーの不等式は次のとおりである。
\[
 \begin{array}{c|ccc}
 \lambda&A_c&A_e&A_p\\ \hline
 5&<1&<1&<-3/16\\
 8&<-1/20&<-2/5&<-3/5.
 \end{array}
\]
さらに、閉領域
\[
 5\leq\lambda\leq8,\qquad(\rho,z)\in\mathcal D,\qquad
 \sqrt{1-z^2}\geq1/\lambda,
\]
において、横断方向の係数は次の上界をもつ。
\[
 \begin{array}{c|c}
 \lambda\text{ interval}&H\text{ upper bound}\\ \hline
 {[5,6]}&-68778309/2^{36}\\
 {[6,7]}&-9210649/2^{33}\\
 {[7,8]}&-1095717/2^{30}.
 \end{array}
\]
\end{computationalassertion}

これらの言明は、有限の積分と有限被覆のみに関するものである。
補集合である極領域を含めた、すべての$\lambda\geq5$への
拡張は、第~\ref{sec:prolate-tail}節における係数の解析的な比較から
従う。特に、$\lambda=8$での数値的な検査を、
その議論なしに外挿しているのではない。
アンカーの完全な厳密端点は、付属の
ソースノートに保持されている\citep{data-chord}。

\begin{computationalassertion}[極での下側包絡積分]
\label{comp:c6}
式~\eqref{eq:pt-cap-computational-integral}で明示的に定義された下界積分
$\mathcal B$に対して
\[
 \mathcal B(\epsilon,w,u)\geq
 \frac{47473218194072764542240096942353646083}{2^{127}}
 >\frac14
 \quad\text{on}\quad
 [0,1/5]\times[0,1]\times[0,1].
\]
\end{computationalassertion}

ここで$z=\sqrt{1-\epsilon^2w^2}$、$\rho=\epsilon wu$である。
$\epsilon>0$に対して、解析的な導出は
$\lambda=1/\epsilon$と同一視し、$\lambda E_z\geq\mathcal B$を証明する。
端点$\epsilon=0$はこの下界積分の連続な
拡張であって、無限大のパラメータにおけるエネルギーの
定義ではない。$\mathcal B$に含まれる解析的な差し引き項は、
省いた角度領域の分を考慮したものである。計算主張~\ref{comp:c6}は
大域的な漸近剰余評価を仮定しない\citep{data-cap}。

\subsection{算術、積分、および特異性の扱い}
\label{subsec:computational-methods}

\paragraph{方向付き丸めによるbinary64包囲。}
有限扁平域、球に近い帯、および扁長側の窓の核は、
IEEE binary64の端点をもつ独自の区間型を用いる。プログラムは
\texttt{FE\_UPWARD}を設定する。算術演算の下端は
適切な符号反転によって得られ、上端は
上向き丸めによって得られる。正の平方根は、上向きの端点と、
その結果を表現可能な値1つ分だけ下げて得られる下端とをもつ。
零は厳密に扱う。記録されたビルドフラグには
\texttt{-frounding-math}、\texttt{-fno-fast-math}、
\texttt{-ffp-contract=off}が含まれる。その正しさは、明示した
浮動小数点環境とコンパイラの挙動に依存するのであって、単に
多くの桁を印字することに依存するのではない。起動時の方向付き丸めの検査が備わっている。

通常の係数表は8,192個のセルと60次の正項
級数をもち、それらの核が用いる係数の量と導関数に対して
外向きの尾の許容量$2^{-38}$をとる。
二進の角度格子は、平方根による半角の
恒等式を通して構成する。外側の箱の座標は
$2^{-20}$を単位とする厳密な整数である。通常報告する積分の端点は
$2^{-36}$を単位とする。扁平側の極での細分については、恒等式
$J'=Q$と$K'=gQ$が、固定した表の節点からの平均値形式の包囲を与える。

空間セル上の主たる積分包囲は、そのセル上での
核全体の値域に、セルの正の測度を掛けたものである。
これらの積の和はDarboux包囲であり、評価されていない
求積剰余を落とすことはない。扁平側の勾配付き座標
$m=1-(1-t)^2$は、$m$におけるセルの厳密な差分を重みとして用いる。
その表示$C=(1-m^2)+\lambda^2m^2$は、薄い極の近くで
正値性を保つ。接触セルでは、零を含む区間で
割る代わりに、既知の幾何学的な比の上下界を用いる。

\paragraph{臨界値の細分のための4次の剰余。}
扁長側の窓における位置決定の計算は、別個のテンソル
Simpson包囲を用いる。その係数表は4,096個のセル、
96次をもち、6階までの正規化した導関数の尾は
$2^{-52}$で抑えられる。幅$\Delta$の小区間上で、Hermite剰余は
符号の定まった次の平均の包囲をもつ。
\[
 \operatorname{mean}(f)-\operatorname{Simpson}(f)
 \in -\frac{\Delta^4}{120}\,[f^{(4)}/4!].
\]
実際、正規化した節点多項式は$[0,1]$上で
$t(t-1)(t-1/2)^2\leq0$であり、その積分は$-1/120$である。
コードは、各小区間全体の上で4次のTaylor係数を包囲する。
2変数の場合、2つの積分誤差は、正の外側の重みを用いて
順に分解する。そのうえで、臨界端点の
包囲を$2^{-47}$を単位として外向きに報告する。
この誤差制御は、通常の点求積や、
2つの数値解像度の間の比較とは別のものである。

\paragraph{Arbによる係数と、有限な裾のための整数核。}
計算主張~\ref{comp:c5}で必要な6つの1次元のアンカーは、
96ビットのArb、$m$の一様な2,048個のセル、および
式~\eqref{eq:pt-positive-series}における$Q$の12次の下側多項式を用いて
包囲する。その省いた係数は
非負であるから、それらを捨てると必要な下界が得られる。
独立な監査は4,096個のセルと14次を用いる。
有限な$[5,8]$の橋は、スケール
$2^{-24}$の固定小数点の区間端点と、符号付き128ビット整数の中間値を用いる。
平方根は、定義となる整数の平方の
不等式が成り立つまで補正する。浮動小数点による初期推定値を、
その補正なしに採用することはない。その2,048セルの係数表は
144ビットのArbによって生成し、16次と、式~\eqref{eq:pt-series-enclosures}における$P$の明示的な正の
尾の上界を用いる。$Q$は
その正の打ち切りによって下から抑える。厳密な三角モーメントは
表の中で包囲する。2つの負の弦の寄与を
対にしたまま保つことで、恒等式
\eqref{eq:pt-paired-negative-identity}が保たれる。

\paragraph{固定された極の被覆。}
計算主張~\ref{comp:c6}は、128ビットのArbと
$16\times8\times8=1,024$個の厳密な外側の箱を用いる。残した
$\theta$の区間は96個の等しいセルに分割し、
$\phi$の区間は48個に分割する。各セル全体での値域に
その面積を掛ける。この残した角度の長方形の上では、分母の正の
下界を計算に先立って証明する。
省いた2つの領域は解析的に抑える。

\paragraph{境界での接触。}
直接の勾配核は、接触点以外のすべての曲面上の点で有界であり、
優関数で抑えられた内側極限をもつ。接触点の面積は零である。
Hessianの包囲では、弦による変数変換が、
消える端点距離を評価の前に相殺する。特に、
つぶれる弦においては有界な弦の表示を拡張する。
一般の特異な表示をそこで数値的に評価することはない。
境界の赤道での導関数にも専用の核がある。
退化した中心に用いる高階の空間微分は、
物理的な境界から離れたところで評価する。

\subsection{どの有限計算が検査されたか}
\label{subsec:computational-workload}
\begin{center}\small
\begin{tabularx}{\linewidth}{lXX}\toprule
入力 & 分類を支える計算 & 追加検査の範囲\\\midrule
C1 & 扁平側の排除と軸上の符号 & 以前に選んだ123項目。扁平側の上側の帯にある277個のうち、正の軸方向勾配の葉で最も弱いものに対する新たなArb検査。\\
C2 & 中心と極の零点、および軸上の根の個数 & $c_o$に対する間隔$10^{-12}$での新たな端点計算。以前の根の個数と単調性の記録が一意性を与える。\\
C3 & 鉛直方向の排除と正の境界トレース & 以前に選んだ35個の葉。28,358個の正の鉛直勾配の葉のうち最も弱いものに対する新たなArb検査。第2の実装が、追加の赤道上の箱176個すべてを包囲し直す。\\
C4 & 鉛直方向の排除、動径方向の単調性、および円の運動 & 動径方向の箱731個すべてを第2の実装で検査した。以前の別個の検査には、鉛直方向の葉22個と、選んだ形状微分の葉1個が含まれる。\\
C5 & 有限の橋と裾の6つのアンカー & 6つのアンカーすべてをArbで再計算した。橋の親79個の細分は、既存の弦の核を用いる。\\
C6 & 極での積分 & 1,024個の外側の箱すべてを、異なる積分実装と角度分割で再計算した。\\\bottomrule
\end{tabularx}\end{center}
補足資料は、元の母集団と、再実行、積分の第2の実装、および別個に導出した公式の間の区別を明らかにする。鉛直方向と形状方向の母集団全体は、動径方向の母集団とは監査の被覆が異なる。
